(1) 线性平面声波的压力与法向质点速度的比值是特征阻抗 \(Z_j=\rho_jc_j>0\)。对于右行声波 \(v=p/Z_j\)，左行波因为传播方向反向而有 \(v=-p/Z_j\)；不能把反射波的压力和速度同时按同一符号叠加。
(2) 设入射压力幅值为 \(p_i\)，反射压力为 \(r_p p_i\)，透射压力为 \(t_pp_i\)。压力连续给 \(1+r_p=t_p\)，速度连续给 \((1-r_p)/Z_1=t_p/Z_2\)。联立求得 \(r_p=(Z_2-Z_1)/(Z_2+Z_1)\)、\(t_p=2Z_2/(Z_1+Z_2)\)；当两介质阻抗不同时，即使没有任何材料耗散也会发生反射。
(3) 无损平面波平均声强大小为 \(I=|p|^2/(2Z)\)，故反射能流率 \(\mathcal R=|r_p|^2=(Z_2-Z_1)^2/(Z_2+Z_1)^2\)，透射能流率 \(\mathcal T=(Z_1/Z_2)|t_p|^2=4Z_1Z_2/(Z_1+Z_2)^2\)。二者严格满足 \(\mathcal R+\mathcal T=1\)，说明较大的透射压力振幅也不必意味着更大的透射功率。
(4) 若 \(Z_1=Z_2\)，则 \(r_p=0,t_p=1\)，无反射。当 \(Z_2\gg Z_1\)，反射压力近似同相 \(r_p\to+1\)，透射压力系数 \(t_p\to2\)，但能量透射率 \(\mathcal T\to0\)；较大压力振幅对应很小的质点速度，因此仍可只有微弱的传输能流。
